\documentclass[11pt]{article}
\usepackage[margin=1in]{geometry}
\usepackage{amsmath,amssymb,amsthm,mathrsfs}
\usepackage{booktabs}
\usepackage{enumitem}
\usepackage[colorlinks=true,linkcolor=blue,citecolor=blue,urlcolor=blue]{hyperref}

\newtheorem{theorem}{Theorem}[section]
\newtheorem{proposition}[theorem]{Proposition}
\newtheorem{lemma}[theorem]{Lemma}
\newtheorem{corollary}[theorem]{Corollary}
\theoremstyle{remark}
\newtheorem{remark}[theorem]{Remark}

\newcommand{\R}{\mathbb{R}}
\newcommand{\T}{\mathbb{T}}
\newcommand{\eps}{\varepsilon}
\newcommand{\dd}{\,\mathrm{d}}
\newcommand{\Lp}{L^p}
\newcommand{\Hs}{H^s}
\newcommand{\Wsp}[2]{W^{#1,#2}}
\newcommand{\norm}[1]{\lVert #1\rVert}
\newcommand{\abs}[1]{\lvert #1\rvert}
\newcommand{\ip}[2]{\langle #1,#2\rangle}
\newcommand{\Div}{\operatorname{div}}
\newcommand{\curl}{\operatorname{curl}}
\newcommand{\Leray}{\mathbb{P}}
\DeclareMathOperator{\sgn}{sgn}

\title{Forced versus unforced Navier--Stokes:\\[2pt]
\large what a prescribed force can and cannot prove}
\author{Working notes prepared with an AI research assistant\\
\small (classical statements are attributed; the arguments of
\S\ref{sec:prescription}--\S\ref{sec:price} are proved in full)}
\date{November 2, 2026}

\begin{document}
\maketitle

\begin{abstract}
This paper does \emph{not} prove the Navier--Stokes Millennium Prize Problem, and no
statement below claims otherwise. Its purpose is to separate, rigorously, three things that
were widely conflated in the 2026 discussion of a claimed AI proof: (i) the \emph{unforced}
problem, which remains open; (ii) the \emph{forced} problem in the formulation of
\cite{Fefferman2000}, which is a control problem and, in a precise sense, trivially solvable;
and (iii) the quantitative price one has to pay, in the size of the force, to make a smooth
flow blow up.

We prove, completely and at the level of detail of a research paper:
\begin{enumerate}[label=(\Alph*),leftmargin=*]
\item the scaling law of the equations and of the force norm (Proposition~\ref{prop:scaling});
\item a prescription theorem: \emph{every} smooth divergence-free path is the velocity field of
      a solution for a suitable smooth force, and that force is unique
      (Theorem~\ref{thm:prescription}); consequently the forced problem has no existence
      content, only a quantitative one;
\item a lower bound for that price: an exactly self-similar blow-up requires
      $\norm{f(t)}_{\infty}\gtrsim (T-t)^{-3/2}$, so a bounded force cannot drive an exactly
      self-similar blow-up (Theorem~\ref{thm:price}, Corollaries~\ref{cor:price} and~\ref{cor:discrete});
\item the conditional regularity criteria of Prodi--Serrin (Theorem~\ref{thm:PS}) and
      Beale--Kato--Majda (Theorem~\ref{thm:BKM}), with complete proofs, and the resulting
      blow-up characterisation;
\item global regularity in two space dimensions (Theorem~\ref{thm:2d}), with a complete proof
      based on the $L^p$ maximum principle for vorticity and on Beale--Kato--Majda.
\end{enumerate}
We then explain, using the quoted theorem of Constantin--Ignatova--Vicol \cite{CIV2026} and its
companion \cite{IKX2026}, why the 2026 construction cannot be converted into a proof for the
unforced problem, and we recall the autoformalisation gap of \cite{BCH2026}: a verified Lean
translation of a paper is not a verification of the paper. Section~\ref{sec:referee} gives a
checklist for refereeing announced proofs of this kind.

Nothing in \S\S\ref{sec:3d}--\ref{sec:2d} is new; those sections are included because a claim
about what a proof must overcome is meaningful only if the surrounding theorems are stated
with their hypotheses intact. The material that is not conveniently available in one place is
\S\ref{sec:prescription}--\S\ref{sec:price} and Appendix~\ref{app:num}.
\end{abstract}

\tableofcontents

\section{Introduction}\label{sec:intro}

\subsection{The problem and its formulations}

Let $\nu>0$ and let $\Omega$ be either $\R^3$ or the flat torus
$\T^3=(\R/2\pi\mathbb{Z})^3$. The incompressible Navier--Stokes equations with a body force
$f$ read
\begin{equation}\label{eq:NS}
\begin{cases}
\partial_t u+(u\cdot\nabla)u-\nu\Delta u+\nabla p=f, & (x,t)\in\Omega\times(0,T),\\
\Div u=0, & (x,t)\in\Omega\times(0,T),\\
u(\cdot,0)=u_0, & x\in\Omega,
\end{cases}
\end{equation}
together with decay at infinity when $\Omega=\R^3$ and with zero spatial mean when
$\Omega=\T^3$.

The Millennium Prize formulation \cite{Fefferman2000} asks, roughly, for a proof or disproof
of the following alternative for $\Omega=\R^3$, $\nu>0$, $f\equiv0$ and smooth, compactly
supported, divergence-free $u_0$: (A) global existence and smoothness of a unique solution;
(B) breakdown, i.e.\ a smooth solution that ceases to be smooth in finite time. The official
statement also permits a prescribed smooth force; the discussion in \cite{SciAm2026} refers to
this variant as option (C). The distinction is not cosmetic. As
Theorem~\ref{thm:prescription} shows, once the force is completely free the equation imposes
no constraint on the \emph{path} of the flow: every smooth divergence-free path is a solution
for a suitable force. What remains is a question about how large, and how regular, that force
must be.

\subsection{The 2026 claims, in the terms used below}

We record the publicly reported facts that motivate this paper, with sources, and no more.
On 8 September 2026 a 166-page manuscript together with a Lean formalisation was announced,
claiming finite-time blow-up for \eqref{eq:NS} in the presence of a smooth external force
\cite{OpenAI2026,WikipediaNS}. On 17 September 2026 Constantin, Ignatova and Vicol posted a
theorem showing that, within the structural class of that construction, a force which is
real-analytic in space (uniformly in time) forces the solution to be regular; in particular
the force in that construction can neither vanish identically near the putative singular
point nor be real-analytic there \cite{CIV2026}. A companion paper replaces spatial
analyticity by rotational invariance of $\curl f$ \cite{IKX2026}. On 6 October 2026 a
collection of 722 manuscripts, grouped into 372 ``result families'', was released; three were
withdrawn the next day because of a sign error, leaving 719, and two further manuscripts were
added on 26 October, leaving 721 manuscripts in 373 families
\cite{OpenAIrepo2026,RW2026}. The Clay Mathematics Institute lists the problem as
\emph{active} and states that its verification is deliberately unhurried \cite{ClayNS}. A
separate analysis argues that the Lean formalisation of the Navier--Stokes manuscript does
not correspond to its natural-language argument \cite{BCH2026}.

Two numbers are frequently conflated in reporting: ``more than 700 problems'' refers to
manuscripts and result families in the October release, whereas the September
Navier--Stokes manuscript is not part of that collection, whose only Navier--Stokes family
concerns universal computation in forced flows \cite{OpenAIrepo2026}. Neither number is a
count of theorems.

\subsection{What is proved here, and what is not}

\begin{center}
\begin{tabular}{@{}ll@{}}
\toprule
Statement & Status in this paper\\
\midrule
3D global regularity / blow-up, $f\equiv0$ & open; \emph{not} claimed\\
2D global regularity & Theorem~\ref{thm:2d} (complete proof)\\
Prodi--Serrin criterion & Theorem~\ref{thm:PS} (complete proof)\\
Beale--Kato--Majda criterion & Theorem~\ref{thm:BKM} (complete proof)\\
Prescription theorem for the forced problem & Theorem~\ref{thm:prescription} (complete proof)\\
$\norm{f(t)}_\infty\gtrsim(T-t)^{-3/2}$ for exact self-similarity
  & Theorem~\ref{thm:price} (complete proof)\\
Consequence of \cite{CIV2026} for the 2026 construction & Corollary~\ref{cor:CIV}
  (deduced from the quoted theorem)\\
Lean verification versus natural-language proof & \S\ref{sec:lean} (quoted, not reproved)\\
\bottomrule
\end{tabular}
\end{center}

The logical status of the paper is therefore the following. We do not know whether the
unforced problem has a negative answer; nobody does. We do know, and prove, exactly why the
forced answer is easy, and why its easiness is paid for by a force that degenerates at the
singular time.

\subsection{Notation}

$L^p=L^p(\Omega)$ and $H^s=H^s(\Omega)$ are the usual Lebesgue and Sobolev spaces;
$L^p_tL^q_x=L^p(0,T;L^q(\Omega))$; $H^s_\sigma$ is the closure in $H^s$ of smooth
divergence-free fields (with zero mean on $\T^3$). $\Leray$ denotes the Leray projector onto
divergence-free fields, $A=-\Leray\Delta$ the Stokes operator, and
$B(u,v)=\Leray\Div(u\otimes v)$. We write $\norm{\cdot}_p=\norm{\cdot}_{L^p}$ and
$\norm{\cdot}=\norm{\cdot}_2$. Constants are denoted by $C$ and may change from line to
line; dependence on parameters is indicated by subscripts. All functions are real-valued.

\section{The equations, their scaling, and their energy structure}\label{sec:scaling}

\subsection{Scaling}

\begin{proposition}[Scaling of the forced equations]\label{prop:scaling}
Let $(u,p,f)$ solve \eqref{eq:NS} on $\Omega\times(0,T)$ with viscosity $\nu>0$ and initial
datum $u_0$. For $\lambda>0$ define
\begin{equation}\label{eq:scaling}
u_\lambda(x,t)=\lambda\,u(\lambda x,\lambda^2 t),\qquad
p_\lambda(x,t)=\lambda^2 p(\lambda x,\lambda^2 t),\qquad
f_\lambda(x,t)=\lambda^3 f(\lambda x,\lambda^2 t),
\end{equation}
for $0<t<T/\lambda^2$. Then $(u_\lambda,p_\lambda,f_\lambda)$ solves \eqref{eq:NS} on
$\Omega\times(0,T/\lambda^2)$ with viscosity $\nu$, force $f_\lambda$ and initial datum
$u_{0,\lambda}(x)=\lambda u_0(\lambda x)$. Conversely, if $(v,q,g)$ solves \eqref{eq:NS} with
viscosity $\nu$ and $(v,q,g)=(u_\lambda,p_\lambda,f_\lambda)$, then $(u,p,f)$ solves
\eqref{eq:NS}.
\end{proposition}

\begin{proof}
Write $y=\lambda x$, $s=\lambda^2 t$ and $U(y,s)=u(y,s)$, so that $u_\lambda(x,t)=\lambda
U(y,s)$ and, by the chain rule,
$\partial_t u_\lambda=\lambda\cdot\lambda^2(\partial_sU)(y,s)=\lambda^3\partial_sU$,
$\nabla_xu_\lambda=\lambda^2(\nabla_yU)(y,s)$ and
$\Delta_xu_\lambda=\lambda^3(\Delta_yU)(y,s)$.
For the convective term, $(u_\lambda\cdot\nabla_x)=\lambda U\cdot\lambda\nabla_y
=\lambda^2(U\cdot\nabla_y)$, hence
$(u_\lambda\cdot\nabla_x)u_\lambda=\lambda^3(U\cdot\nabla_y)U$. Similarly
$\nabla_xp_\lambda=\lambda^3\nabla_yP$ and $\Div_xu_\lambda=\lambda^2(\Div_yU)=0$. Dividing
the momentum equation by $\lambda^3$ gives the equation at $(y,s)$. The converse is the same
computation with $\lambda$ replaced by $\lambda^{-1}$.
\end{proof}

The action of the scaling on norms is what makes the problem critical. From
\eqref{eq:scaling}, for $1\le q\le\infty$,
\begin{equation}\label{eq:normscaling}
\norm{u_\lambda(t)}_{q}=\lambda^{1-3/q}\norm{u(\lambda^2t)}_{q},\qquad
\norm{\nabla u_\lambda(t)}_{2}=\lambda^{1/2}\norm{\nabla u(\lambda^2t)}_{2},
\end{equation}
and
\begin{equation}\label{eq:forcescaling}
\norm{f_\lambda(t)}_{q}=\lambda^{3-3/q}\norm{f(\lambda^2t)}_{q},\qquad
\norm{f_\lambda}_{L^1(0,T/\lambda^2;L^\infty)}=\lambda\,\norm{f}_{L^1(0,T;L^\infty)}.
\end{equation}
Thus $\norm{u}_{3}$ and $\norm{u}_{H^{1/2}}$ are invariant, $\norm{\nabla u}_{2}$ grows like
$\lambda^{1/2}$, and a bounded force is \emph{not} small in scaling terms: its $L^\infty$ norm
grows like $\lambda^{3}$. This single computation already explains why statements about the
forced problem are statements about the size of $f$, not about the size of $u$.

\subsection{Energy identity and its consequences}

\begin{proposition}[Energy identity]\label{prop:energy}
Let $u$ be a smooth solution of \eqref{eq:NS} on $\T^3\times[0,T]$, or a Schwartz-class
solution on $\R^3\times[0,T]$. Then for $0\le t\le T$,
\begin{equation}\label{eq:energy}
\frac12\frac{\dd}{\dd t}\norm{u(t)}^2+\nu\norm{\nabla u(t)}^2=\ip{f(t)}{u(t)}.
\end{equation}
Consequently
\begin{equation}\label{eq:energybound}
\norm{u(t)}^2\le e^{\nu t}\norm{u_0}^2+\frac1\nu\int_0^te^{\nu(t-s)}\norm{f(s)}^2\dd s .
\end{equation}
\end{proposition}

\begin{proof}
Take the $L^2$ scalar product of the momentum equation with $u$ and integrate by parts. The
pressure term vanishes because $\Div u=0$ (on $\T^3$ by periodicity, on $\R^3$ by decay), and
the convective term vanishes because
\[
\int_\Omega (u\cdot\nabla)u\cdot u\dd x
=\frac12\int_\Omega u\cdot\nabla\abs{u}^2\dd x
=-\frac12\int_\Omega(\Div u)\abs{u}^2\dd x=0 .
\]
The viscous term gives $\nu\norm{\nabla u}^2$. This is \eqref{eq:energy}. For
\eqref{eq:energybound}, estimate $\ip{f}{u}\le\norm{f}\norm{u}$ and use
$2\norm{f}\norm{u}\le\nu^{-1}\norm{f}^2+\nu\norm{u}^2$, so that with $y=\norm{u}^2$,
$y'\le\nu y+\nu^{-1}\norm{f}^2$; Gronwall's lemma gives \eqref{eq:energybound}.
\end{proof}

\begin{remark}[What the energy controls]\label{rem:energy}
Estimate \eqref{eq:energybound} gives, for $f\in L^2(0,T;L^2)$,
\[
u\in L^\infty(0,T;L^2)\cap L^2(0,T;H^1),
\]
with bounds depending only on $\norm{u_0}$, $\norm{f}_{L^2_tL^2_x}$ and $T$. It gives no
bound on $\norm{\nabla u}_\infty$, on $\norm{u}_\infty$, or on any critical norm. In three
dimensions all three are exactly what the known continuation criteria require
(\S\ref{sec:3d}), and none of them is available from the energy alone. In two dimensions the
situation is different, because the vorticity equation is a transport equation with no
stretching term (\S\ref{sec:2d}).
\end{remark}

\section{Weak solutions: what Leray and Hopf already give}\label{sec:weak}

\begin{theorem}[Leray \cite{Leray1934}, Hopf \cite{Hopf1951}]\label{thm:lerayhopf}
Let $\Omega\in\{\R^3,\T^3\}$, $\nu>0$, $u_0\in L^2_\sigma(\Omega)$ and
$f\in L^2(0,T;H^{-1}(\Omega))$. Then there exists
\[
u\in L^\infty(0,T;L^2_\sigma)\cap L^2(0,T;H^1_\sigma)
\]
which satisfies \eqref{eq:NS} in the sense of distributions and the energy inequality
\begin{equation}\label{eq:energyineq}
\frac12\norm{u(t)}^2+\nu\int_0^t\norm{\nabla u(s)}^2\dd s
\le\frac12\norm{u_0}^2+\int_0^t\ip{f(s)}{u(s)}\dd s
\end{equation}
for almost every $t\in[0,T]$ (and for every $t$ after a modification on a null set).
\end{theorem}

\begin{proof}[Proof, in the form needed later]
We recall the standard Galerkin argument because the uniform estimates are used in
\S\ref{sec:2d} and \S\ref{sec:3d}.

\emph{Step 1 (Galerkin system).} Let $(\varphi_k)_{k\ge1}$ be the eigenfunctions of the
Stokes operator $A$ on $\Omega$, orthonormal in $L^2_\sigma$ with eigenvalues
$0<\lambda_1\le\lambda_2\le\cdots\to\infty$ (on $\T^3$ restricted to zero mean). Let
$V_m=\operatorname{span}\{\varphi_1,\dots,\varphi_m\}$ and look for
$u_m(t)=\sum_{k\le m}c_{km}(t)\varphi_k$ solving
\[
\ip{\partial_tu_m}{\varphi}+\nu\ip{\nabla u_m}{\nabla\varphi}
+\ip{(u_m\cdot\nabla)u_m}{\varphi}=\ip{f}{\varphi},
\qquad \varphi\in V_m,
\]
with $u_m(0)=\Leray_m u_0$. This is a quadratic ODE system with smooth coefficients, hence
has a local solution; the estimate below makes it global in $[0,T]$.

\emph{Step 2 (uniform energy estimate).} Taking $\varphi=u_m$ and using the cancellation of
Proposition~\ref{prop:energy},
\[
\frac12\frac{\dd}{\dd t}\norm{u_m}^2+\nu\norm{\nabla u_m}^2=\ip{f}{u_m}
\le\norm{f}_{H^{-1}}\norm{\nabla u_m},
\]
hence, after absorbing and integrating,
\begin{equation}\label{eq:galerkin}
\sup_{0\le t\le T}\norm{u_m(t)}^2+\nu\int_0^T\norm{\nabla u_m}^2\dd t
\le C\bigl(\norm{u_0}^2+\norm{f}_{L^2(0,T;H^{-1})}^2\bigr),
\end{equation}
with $C$ independent of $m$.

\emph{Step 3 (compactness).} By \eqref{eq:galerkin}, $(u_m)$ is bounded in
$L^\infty(0,T;L^2_\sigma)\cap L^2(0,T;H^1_\sigma)$, and $\partial_tu_m$ is bounded in
$L^2(0,T;H^{-2})$ (bound the three terms of the equation separately, using
$\norm{B(u_m,u_m)}_{H^{-2}}\le C\norm{u_m}\norm{\nabla u_m}$). By the Aubin--Lions lemma,
after a subsequence $u_m\to u$ strongly in $L^2(0,T;L^2_{\mathrm{loc}})$ and weakly in the
two spaces above; the limit is divergence-free and satisfies the equation.

\emph{Step 4 (energy inequality).} The Galerkin identity gives, for every $m$ and a.e.\ $t$,
\[
\frac12\norm{u_m(t)}^2+\nu\int_0^t\norm{\nabla u_m}^2
=\frac12\norm{u_m(0)}^2+\int_0^t\ip{f}{u_m}.
\]
The norm is weakly lower semicontinuous and $\ip{f}{u_m}\to\ip{f}{u}$ because
$u_m\to u$ strongly in $L^2_tL^2_x$; passing to the limit gives \eqref{eq:energyineq}.
\end{proof}

\begin{remark}\label{rem:weak}
Two points matter for the rest of the paper. First, Theorem~\ref{thm:lerayhopf} is an
\emph{existence} theorem: it produces a solution for every divergence-free $L^2$ datum and
every $f\in L^2_tH^{-1}_x$. It says nothing about uniqueness or smoothness. Second, in two
dimensions the weak solution is unique and, as Theorem~\ref{thm:2d} shows, smooth for $t>0$;
in three dimensions both properties are open, and the Millennium Problem is precisely the
question whether smoothness can fail.
\end{remark}

\section{Conditional regularity in three dimensions}\label{sec:3d}

The results of this section are classical. They are stated with full hypotheses and proved,
because they are exactly the statements that a claimed proof of the unforced problem must
either establish or circumvent.

\begin{lemma}[Local existence and continuation]\label{lem:local}
Let $\Omega\in\{\R^3,\T^3\}$, $s>1/2$, $u_0\in H^s_\sigma$ and
$f\in L^2(0,T;H^{s-1})$. There exist $T_0=T_0(\norm{u_0}_{H^s})>0$ and a unique
\[
u\in C([0,T_0];H^s_\sigma)\cap L^2(0,T_0;H^{s+1}),\qquad
\partial_tu\in L^2(0,T_0;H^{s-1}),
\]
solving \eqref{eq:NS}. If $T^*$ is the maximal time of existence, then either $T^*=\infty$ or
\begin{equation}\label{eq:blowupcrit}
\limsup_{t\to T^*}\norm{u(t)}_{H^s}=\infty .
\end{equation}
Moreover, for $s=1$ the lifespan can be chosen as
$T_0=c\,(1+\norm{\nabla u_0})^{-4}$ with $c>0$ absolute, and the solution depends
continuously on the data.
\end{lemma}

\begin{proof}[Reference]
This is Kato's theory \cite{Kato1984}; see also \cite[Chapter 3]{RRST2001} and
\cite[Chapter 5]{Lemarie2016}. The key point for us is the uniformity of $T_0$: it depends
only on the $H^s$ norm of the data, so that a uniform $H^s$ bound on $[0,T^*)$ permits
restarting at $T^*$ and iterating. We use this in the proofs of
Theorems~\ref{thm:PS} and~\ref{thm:BKM}.
\end{proof}

\subsection{The Prodi--Serrin criterion}

\begin{theorem}[Prodi \cite{Prodi1959}, Serrin \cite{Serrin1962}]\label{thm:PS}
Let $u$ be the smooth solution of \eqref{eq:NS} on $[0,T^*)$ with
$f\in L^2(0,T^*;L^2)$ and let $q>3$, $p\in[2,\infty]$ satisfy
\begin{equation}\label{eq:PScond}
\frac2p+\frac3q\le1 .
\end{equation}
If $u\in L^p(0,T^*;L^q)$, then $T^*=\infty$ (or, if the solution is considered on a bounded
interval, it extends beyond $T^*$).
\end{theorem}

\begin{proof}
Let $y(t)=\norm{\nabla u(t)}^2$ for $t<T^*$. We test the momentum equation with $-\Delta u$:
\[
\frac12y'(t)+\nu\norm{\Delta u}^2
=\int_\Omega(u\cdot\nabla)u\cdot\Delta u\dd x-\int_\Omega f\cdot\Delta u\dd x .
\]
(The pressure term disappears because $\Div u=0$.) We bound the two terms separately.

\emph{The convective term.} By H\"older's inequality with exponents $q$,
$r=2q/(q-2)$ and $2$,
\[
\abs{\int(u\cdot\nabla)u\cdot\Delta u}
\le\norm{u}_{q}\norm{\nabla u}_{r}\norm{\Delta u}_{2},
\]
which is legitimate because $1/q+1/r+1/2=1$. The Gagliardo--Nirenberg inequality in
$\R^3$ (Appendix~\ref{app:ineq}) gives, with $\theta=3/q\in(0,1)$,
\[
\norm{\nabla u}_{r}\le C\norm{\nabla u}_{2}^{1-\theta}\norm{\Delta u}_{2}^{\theta}.
\]
Hence
\[
\abs{\int(u\cdot\nabla)u\cdot\Delta u}
\le C\norm{u}_{q}\norm{\nabla u}_{2}^{1-3/q}\norm{\Delta u}_{2}^{1+3/q}.
\]
By Young's inequality $ab\le\eps b^{\sigma}+C_\eps a^{\sigma'}$ with
$\sigma=2/(1+3/q)=2q/(q+3)$ and $\sigma'=2q/(q-3)$,
\[
\abs{\int(u\cdot\nabla)u\cdot\Delta u}
\le\frac{\nu}{2}\norm{\Delta u}^2
+C_\nu\norm{u}_{q}^{2q/(q-3)}\norm{\nabla u}^2 .
\]

\emph{The force term.} $\abs{\int f\cdot\Delta u}\le\frac{\nu}{4}\norm{\Delta u}^2
+\frac1\nu\norm{f}^2$.

Combining, and discarding the non-negative term $\frac{\nu}{4}\norm{\Delta u}^2$,
\begin{equation}\label{eq:PSode}
y'(t)\le C_\nu\norm{u(t)}_{q}^{2q/(q-3)}y(t)+\frac{2}{\nu}\norm{f(t)}^2 .
\end{equation}
Gronwall's lemma gives, for $t<T^*$,
\begin{equation}\label{eq:PSgronwall}
y(t)\le\Bigl(y(0)+\frac2\nu\int_0^t\norm{f}^2\Bigr)
\exp\Bigl(C_\nu\int_0^t\norm{u(s)}_{q}^{2q/(q-3)}\dd s\Bigr).
\end{equation}
It remains to see that the exponent is finite. Put $s_0=2q/(q-3)$. Condition
\eqref{eq:PScond} is equivalent to $s_0\le p$, and H\"older's inequality in time gives
\[
\int_0^{T^*}\norm{u}_{q}^{s_0}\dd s
\le(T^*)^{1-s_0/p}\Bigl(\int_0^{T^*}\norm{u}_{q}^{p}\dd s\Bigr)^{s_0/p}
\le(T^*)^{1-s_0/p}\norm{u}_{L^p(0,T^*;L^q)}^{s_0}<\infty .
\]
Therefore $\sup_{0\le t<T^*}\norm{\nabla u(t)}\le M<\infty$. By
Lemma~\ref{lem:local}, the $H^1$ solution can be restarted at any time $t_0<T^*$ with a
lifespan bounded below by $c(1+M)^{-4}$, so finitely many steps cover $[0,T^*]$ and extend
the solution beyond $T^*$. Hence $T^*=\infty$.
\end{proof}

\begin{remark}[The critical case $q=3$]\label{rem:q3}
At $q=3$ the exponent $s_0=2q/(q-3)$ is infinite and the argument breaks down; the
corresponding statement (smallness of $\norm{u}_{L^\infty_tL^3}$, or of
$\norm{u}_{L^3}$ globally) is the borderline case and is open for large data. This is the
analytic expression of the fact that $L^3$ is the scaling-invariant Lebesgue space
(\S\ref{sec:scaling}).
\end{remark}

\subsection{A priori $L^q$ bounds}

\begin{corollary}[An a priori $L^q$ bound suffices]\label{cor:subcritical}
Let $q>3$ and suppose that the smooth solution of \eqref{eq:NS} on $[0,T^*)$ satisfies
\begin{equation}\label{eq:Lqbound}
\sup_{0\le t<T^*}\norm{u(t)}_{q}=:M_q<\infty .
\end{equation}
Then $T^*=\infty$.
\end{corollary}

\begin{proof}
Apply Theorem~\ref{thm:PS} with $p=\infty$: condition \eqref{eq:PScond} reads $3/q\le1$,
which holds strictly because $q>3$, and $u\in L^\infty(0,T^*;L^q)$ by \eqref{eq:Lqbound}.
\end{proof}

\begin{remark}\label{rem:subcritical}
Corollary~\ref{cor:subcritical} is the quantitative content of the word ``subcritical'': a
uniform bound in any $L^q$ with $q>3$ rules out blow-up. The corollary is, however, only a
conditional statement, and it is worth being explicit about why the bound \eqref{eq:Lqbound}
is not available for free. Testing the equation with $u\abs{u}^{q-2}$ gives
\[
\frac1q\frac{\dd}{\dd t}\norm{u}_{q}^{q}
=\nu\int_\Omega\abs{u}^{q-2}u\cdot\Delta u\dd x
-\int_\Omega\abs{u}^{q-2}u\cdot\nabla p\dd x ,
\]
the convective term integrating to zero because $\Div u=0$. The viscous term is non-positive
(Appendix~\ref{app:ineq}), but the pressure term equals
$\int_\Omega p\,\Div(\abs{u}^{q-2}u)=\int_\Omega p\,u\cdot\nabla\abs{u}^{q-2}$, which is a
nonlinear term with no sign and does not vanish in general. In particular, unlike the heat equation, the
Navier--Stokes flow is not known to decrease its $L^q$ norms for $q>2$, and no bound of the
form \eqref{eq:Lqbound} follows from the energy inequality. Obtaining such a bound is
equivalent to a substantial part of the problem.
\end{remark}


\subsection{The Beale--Kato--Majda criterion}

\begin{theorem}[Beale--Kato--Majda \cite{BKM1984}]\label{thm:BKM}
Let $\Omega=\R^3$ or $\T^3$, $s>5/2$, and let $u$ be the smooth solution of \eqref{eq:NS} on
$[0,T^*)$ with $f\in L^1(0,T^*;H^s)$. If
\begin{equation}\label{eq:BKMcond}
\int_0^{T^*}\norm{\nabla u(t)}_{\infty}\dd t=:L<\infty,
\end{equation}
then $T^*=\infty$.
\end{theorem}

\begin{proof}
Apply the fractional derivative $\Lambda^s=(-\Delta)^{s/2}$ to the equation, take the $L^2$
scalar product with $\Lambda^su$, and use that $\Lambda^s\nabla p$ is orthogonal to
$\Lambda^su$ (because $\Div u=0$):
\[
\frac12\frac{\dd}{\dd t}\norm{\Lambda^su}^2+\nu\norm{\Lambda^s\nabla u}^2
=-\ip{[\Lambda^s,u\cdot\nabla]u}{\Lambda^su}+\ip{\Lambda^sf}{\Lambda^su}.
\]
The Kato--Ponce commutator estimate (Appendix~\ref{app:ineq}) gives, for $s>5/2$,
\[
\abs{\ip{[\Lambda^s,u\cdot\nabla]u}{\Lambda^su}}
\le C\norm{\nabla u}_{\infty}\norm{\Lambda^su}^2 .
\]
Hence, with $Y(t)=\norm{u(t)}_{H^s}$,
\[
Y'(t)\le C\norm{\nabla u(t)}_{\infty}Y(t)+\norm{f(t)}_{H^s}.
\]
Gronwall's lemma yields
\[
\sup_{0\le t<T^*}\norm{u(t)}_{H^s}
\le\Bigl(\norm{u_0}_{H^s}+\int_0^{T^*}\norm{f}_{H^s}\dd t\Bigr)e^{CL}<\infty .
\]
By the continuation part of Lemma~\ref{lem:local} (with $H^s$ local existence and lifespan
depending only on the $H^s$ norm), the solution extends beyond $T^*$.
\end{proof}

\begin{corollary}[Structure of a blow-up]\label{cor:blowup}
Let $T^*<\infty$ be the maximal time of the smooth solution of \eqref{eq:NS} with
$f\in L^1(0,T^*;H^s)$, $s>5/2$. Then
\begin{equation}\label{eq:blowup}
\int_0^{T^*}\norm{\nabla u(t)}_{\infty}\dd t=\infty,
\qquad
\limsup_{t\to T^*}\norm{\nabla u(t)}_{\infty}=\infty,
\qquad
\limsup_{t\to T^*}\norm{u(t)}_{H^s}=\infty .
\end{equation}
\end{corollary}

\begin{proof}
The first two assertions are the contrapositive of Theorem~\ref{thm:BKM}; the third is
\eqref{eq:blowupcrit}.
\end{proof}

\begin{remark}\label{rem:BKMmeaning}
Corollary~\ref{cor:blowup} is the precise statement of what a finite-time singularity would
have to look like: the gradient must become unbounded, and not merely integrable, in
$L^\infty$; equivalently, the solution must leave every subcritical Sobolev space. It also
shows why a blow-up construction has to fight the scaling: to produce
$\norm{\nabla u(t)}_\infty\to\infty$ one needs a mechanism that defeats the diffusive
smoothing at scales comparable to the Kolmogorov length, and no such mechanism is known for
$f\equiv0$.
\end{remark}

\section{Two dimensions: a complete regularity proof}\label{sec:2d}

The two-dimensional case is the one place where the programme of
\S\ref{sec:3d} succeeds completely, and the reason is structural: in two dimensions the
vorticity equation has no stretching term, and the vorticity is transported and diffused
rather than amplified. We give the full argument because it shows exactly which quantity
replaces $\norm{\nabla u}_\infty$ in the continuation criterion, and why that quantity fails
to be controlled in three dimensions.

Throughout this section $\Omega=\T^2$, $\nu>0$, $u_0\in H^1_\sigma(\T^2)$ and
$f\in C^\infty(\T^2\times[0,\infty))$. We write
\[
\omega=\partial_1u_2-\partial_2u_1,\qquad g=\partial_1f_2-\partial_2f_1 .
\]

\begin{theorem}[Global regularity in 2D; Ladyzhenskaya {\rm\cite{Ladyzhenskaya1969}}]
\label{thm:2d}
Under the above hypotheses there is a unique solution
\[
u\in C([0,\infty);H^1_\sigma(\T^2))\cap C^\infty(\T^2\times(0,\infty))
\]
of \eqref{eq:NS}. It satisfies the energy identity \eqref{eq:energy} for all $t\ge0$, and
\[
\sup_{t\in[0,T]}\norm{\nabla u(t)}_{p}<\infty
\qquad\text{for every }T<\infty,\ p\in[2,\infty).
\]
\end{theorem}

\begin{lemma}[Vorticity: transport, diffusion, and $L^p$ bounds]\label{lem:vorticity}
Let $u$ be a smooth solution on $[0,T^*)$. Then $\omega$ solves
\begin{equation}\label{eq:vorticity}
\partial_t\omega+(u\cdot\nabla)\omega=\nu\Delta\omega+g,
\end{equation}
and for every $p\in[2,\infty]$,
\begin{equation}\label{eq:vortbound}
\norm{\omega(t)}_{p}\le\norm{\omega_0}_{p}+\int_0^t\norm{g(s)}_{p}\dd s,
\qquad 0\le t<T^*.
\end{equation}
\end{lemma}

\begin{proof}
Equation \eqref{eq:vorticity} follows by applying $\curl$ to the momentum equation and using
$\curl((u\cdot\nabla)u)=u\cdot\nabla\omega$ for divergence-free $u$ in two dimensions, and
$\curl\nabla p=0$.

For $2\le p<\infty$, multiply \eqref{eq:vorticity} by $\omega\abs{\omega}^{p-2}$ and
integrate over $\T^2$:
\[
\frac1p\frac{\dd}{\dd t}\norm{\omega}_{p}^{p}
=\nu\int_{\T^2}\Delta\omega\,\omega\abs{\omega}^{p-2}\dd x
+\int_{\T^2}g\,\omega\abs{\omega}^{p-2}\dd x,
\]
the transport term integrating to zero because $\Div u=0$ and $\nabla\abs{\omega}^{p}=p\abs{\omega}^{p-2}\omega\nabla\omega$. For the
diffusion term we use Kato's inequality
$\Delta\abs{\omega}^{p}\ge p\abs{\omega}^{p-2}\omega\Delta\omega$, valid for $p\ge2$ in the
sense of distributions (Appendix~\ref{app:ineq}); since $\int_{\T^2}\Delta\abs{\omega}^{p}=0$,
\[
\nu\int_{\T^2}\Delta\omega\,\omega\abs{\omega}^{p-2}\dd x
\le\frac{\nu}{p}\int_{\T^2}\Delta\abs{\omega}^{p}\dd x=0 .
\]
The forcing term is bounded by $\norm{g}_{p}\norm{\omega}_{p}^{p-1}$. Hence
$\frac{\dd}{\dd t}\norm{\omega}_{p}\le\norm{g}_{p}$, which integrates to \eqref{eq:vortbound}
for $p<\infty$. For $p=\infty$ we use
$\norm{\omega}_{p}\uparrow\norm{\omega}_{\infty}$ and the monotone convergence theorem
applied to $\norm{g}_{p}\uparrow\norm{g}_{\infty}$ (both functions are in
$L^\infty$ on the torus), and pass to the limit in \eqref{eq:vortbound}.
\end{proof}

\begin{lemma}[Biot--Savart and Sobolev in 2D]\label{lem:biot}
Let $\omega\in L^p(\T^2)$ with zero mean and let $u=\nabla^\perp(-\Delta)^{-1}\omega$, i.e.
$u=(-\partial_2,\partial_1)(-\Delta)^{-1}\omega$. Then $\curl u=\omega$, $\Div u=0$, and for
every $p\in(1,\infty)$,
\begin{equation}\label{eq:biot}
\norm{\nabla u}_{p}\le C_p\norm{\omega}_{p}.
\end{equation}
Moreover $W^{1,p}(\T^2)\subset L^\infty(\T^2)$ for $p>2$, and
$W^{2,p}(\T^2)\subset W^{1,\infty}(\T^2)$ for $p>2$.
\end{lemma}

\begin{proof}
The operator $\nabla\nabla^\perp(-\Delta)^{-1}$ has symbol
$-\xi\xi^\perp/\abs{\xi}^2$, a matrix of Riesz transforms, which is bounded on $L^p$ for
$1<p<\infty$ (Appendix~\ref{app:ineq}). The Sobolev embeddings in two dimensions are
standard: $W^{1,p}\subset L^\infty$ for $p>2$ and $W^{2,p}\subset C^{1,1-2/p}\subset
W^{1,\infty}$ for $p>2$.
\end{proof}

\begin{lemma}[Maximal regularity for the Stokes system]\label{lem:maxreg}
Let $p\in(1,\infty)$, $F\in L^\infty(0,T;L^p(\T^2))$, $u_0\in L^p_\sigma(\T^2)$, and let
$u$ solve $\partial_tu-\nu\Delta u+\nabla p=F$, $\Div u=0$, $u(0)=u_0$. Then for every
$\delta\in(0,T)$,
\begin{equation}\label{eq:maxreg}
\sup_{\delta\le t\le T}\bigl(\norm{u(t)}_{W^{2,p}}+\norm{\partial_tu(t)}_{p}\bigr)
\le C(\delta,p,\nu,T)\bigl(\norm{u_0}_{p}+\norm{F}_{L^\infty(0,T;L^p)}\bigr).
\end{equation}
\end{lemma}

\begin{proof}[Reference]
This is the maximal $L^p$ regularity of the Stokes semigroup on a torus
\cite{GigaSohr1991}; the constant blows up as $\delta\downarrow0$ like $\delta^{-1}$, which is
harmless for us because the solution is smooth on $[0,\delta/2]$ by
Lemma~\ref{lem:local}.
\end{proof}

\begin{proof}[Proof of Theorem~\ref{thm:2d}]
\emph{Step 1 (local existence).} By Lemma~\ref{lem:local} with $s=1$ there is a unique smooth
solution on a maximal interval $[0,T^*)$.

\emph{Step 2 (uniform vorticity bounds).} By Lemma~\ref{lem:vorticity}, for every
$p\in[2,\infty]$ and $0\le t<T^*$,
\[
\norm{\omega(t)}_{p}\le\norm{\omega_0}_{p}+\int_0^{t}\norm{g(s)}_{p}\dd s
\le\norm{\omega_0}_{p}+T^*\norm{g}_{L^\infty(0,T^*;L^p)}=:M_p .
\]
In particular $\omega\in L^\infty(0,T^*;L^p)$ for every $p<\infty$.

\emph{Step 3 (velocity bounds).} Fix $p>2$. By Lemma~\ref{lem:biot},
$\nabla u\in L^\infty(0,T^*;L^p)$ with norm at most $C_pM_p$, hence
$u\in L^\infty(0,T^*;L^\infty)$ and therefore
\[
B(u,u)=\Leray\Div(u\otimes u)\in L^\infty(0,T^*;L^p),
\qquad
\norm{B(u,u)}_{p}\le C\norm{u}_{\infty}\norm{\nabla u}_{p}.
\]

\emph{Step 4 (bootstrap to a gradient bound).} Apply Lemma~\ref{lem:maxreg} to the equation
$\partial_tu-\nu\Delta u+\nabla p=-B(u,u)+f$ with $F=-B(u,u)+f\in L^\infty(0,T^*;L^p)$. For
every $\delta\in(0,T^*)$,
\[
\sup_{\delta\le t<T^*}\norm{u(t)}_{W^{2,p}}<\infty,
\]
so by Lemma~\ref{lem:biot}, $\nabla u\in L^\infty(\delta,T^*;L^\infty)$. Since $u$ is smooth
on $[0,\delta]$, we conclude
\[
\int_0^{T^*}\norm{\nabla u(t)}_{\infty}\dd t<\infty .
\]

\emph{Step 5 (continuation).} Theorem~\ref{thm:BKM} applies in two dimensions as well (the
commutator estimate is dimension-free and $s>1+2/2=2$ suffices), so $T^*=\infty$. The energy
identity follows from smoothness by Proposition~\ref{prop:energy}, and the bounds on
$\norm{\nabla u}_p$ follow from Steps 2--3.
\end{proof}

\begin{proposition}[Uniqueness in 2D]\label{prop:uniqueness2d}
Let $u,v$ be two Leray--Hopf solutions of \eqref{eq:NS} on $\T^2$ with the same data in
$L^2_\sigma$ and the same $f\in L^2(0,T;L^2)$. Then $u=v$.
\end{proposition}

\begin{proof}
Let $w=u-v$ and $q$ the corresponding pressure. Then
$\partial_tw+(u\cdot\nabla)w+(w\cdot\nabla)v-\nu\Delta w+\nabla q=0$ and $\Div w=0$. Testing
with $w$,
\[
\frac12\frac{\dd}{\dd t}\norm{w}^2+\nu\norm{\nabla w}^2
=-\int_{\T^2}(w\cdot\nabla)v\cdot w\dd x
\le\norm{w}_{4}^2\norm{\nabla v}_2
\le C\norm{w}\norm{\nabla w}\norm{\nabla v}_2,
\]
where we used the two-dimensional Ladyzhenskaya inequality
$\norm{w}_4^2\le C\norm{w}\norm{\nabla w}$ (Appendix~\ref{app:ineq}). Hence
\[
\frac{\dd}{\dd t}\norm{w}^2+\nu\norm{\nabla w}^2\le\frac{C}{\nu}\norm{\nabla v}_2^2\norm{w}^2 .
\]
Since $\norm{\nabla v}_2\in L^2(0,T)$, Gronwall's lemma gives $w=0$.
\end{proof}

\begin{remark}[Where the argument stops in 3D]\label{rem:2d3d}
Three ingredients are used in the proof of Theorem~\ref{thm:2d} and fail in 3D.
(i) The vorticity equation has no stretching term, so \eqref{eq:vortbound} is linear; in 3D
the term $(\omega\cdot\nabla)u$ is present and produces the vortex-stretching term
$\int(\omega\cdot\nabla)u\cdot\omega$, which has no sign.
(ii) The two-dimensional Sobolev embeddings in Lemma~\ref{lem:biot} give $L^\infty$ from $L^p$
gradients at fixed $p>2$; in 3D, $W^{1,p}\subset L^\infty$ requires $p>3$, and the vorticity
bound does not deliver $p>3$ from the energy.
(iii) In 3D the analogue of \eqref{eq:vortbound} is the enstrophy balance, whose right-hand
side contains the vortex-stretching term $\int(\omega\cdot\nabla)u\cdot\omega$; there is no
sign, and no $L^p$ bound for $p>3$ follows. Since $W^{1,p}\subset L^\infty$ in 3D needs
$p>3$, the 2D bootstrap has no counterpart.
None of these obstructions is a proof of anything; they merely explain why the 2D argument
cannot be transplanted.
\end{remark}

\section{Partial regularity: what is known unconditionally}\label{sec:partial}

The strongest unconditional statement in 3D is not regularity but \emph{partial} regularity.

\begin{theorem}[Caffarelli--Kohn--Nirenberg {\rm\cite{CKN1982}}]\label{thm:CKN}
Let $u$ be a suitable weak solution of \eqref{eq:NS} on $\R^3\times(0,T)$ with
$f\in L^2(0,T;L^2_{\mathrm{loc}})$. Then the set $\Sigma\subset\R^3\times(0,T)$ of points in
whose neighbourhood $u$ is unbounded has zero one-dimensional parabolic Hausdorff measure:
$\mathcal{P}^1(\Sigma)=0$. In particular, for almost every $t\in(0,T)$ the set
$\{x:(x,t)\in\Sigma\}$ has zero two-dimensional Hausdorff measure.
\end{theorem}

The proof rests on an $\eps$-regularity criterion: there is an absolute $\eps_*>0$ such that a
suitable weak solution satisfying
\[
\limsup_{r\to0}\frac1r\int_{Q_r(x_0,t_0)}\abs{\nabla u}^2\dd x\dd t<\eps_*
\]
is H\"older continuous in a neighbourhood of $(x_0,t_0)$; here $Q_r$ is the parabolic
cylinder of radius $r$. Theorem~\ref{thm:CKN} then follows from a covering argument.

\begin{remark}[What partial regularity does and does not give]\label{rem:CKN}
Theorem~\ref{thm:CKN} is compatible with both answers to the Millennium Problem: it says that
a singular set, if non-empty, is small, not that it is empty. It gives no bound on the number
of singular points, no stability, and no information at the critical scale. It cannot be used
to rule out the blow-up announced in 2026, nor to support it: a single point $(x_*,T)$ is a set
of parabolic measure zero.
\end{remark}

\section{What a free force can do}\label{sec:prescription}

We now come to the structural reason why the forced problem is not the unforced problem. The
observation is elementary, and it is why the two must be discussed separately.

\begin{theorem}[Prescription theorem]\label{thm:prescription}
Let $T\in(0,\infty]$, $\Omega\in\{\R^3,\T^3\}$, $\nu>0$, and let
$v\in C^\infty(\Omega\times[0,T))$ be divergence-free with $v(\cdot,0)=u_0$. Let $q$ be the
unique zero-mean solution of
\begin{equation}\label{eq:pressure}
-\Delta q=\Div\bigl((v\cdot\nabla)v\bigr)=\partial_i\partial_j(v_iv_j),
\end{equation}
and define
\begin{equation}\label{eq:force}
f:=\partial_tv+(v\cdot\nabla)v-\nu\Delta v+\nabla q .
\end{equation}
Then $f\in C^\infty(\Omega\times[0,T))$, $\Div f=0$, and $(v,q,f)$ solves \eqref{eq:NS} on
$\Omega\times[0,T)$ with initial datum $u_0$. Conversely, every classical solution
$(u,p,f)$ of \eqref{eq:NS} on $[0,T)$ arises in this way from $v=u$, $q=p$.
\end{theorem}

\begin{proof}
Since $\Div v=0$, we have $\Div((v\cdot\nabla)v)=\partial_i\partial_j(v_iv_j)$; indeed
$\Div(v\otimes v)=(v\cdot\nabla)v+(\Div v)v=(v\cdot\nabla)v$. Substituting
\eqref{eq:force} into the momentum equation shows that the left-hand side equals $f$. For the
divergence, $\Div(\partial_tv)=0$ and $\Div\Delta v=\Delta\Div v=0$, hence
\[
\Div f=\Div\bigl((v\cdot\nabla)v\bigr)+\Delta q=0
\]
by \eqref{eq:pressure}. The converse is immediate from the definition of a classical
solution.
\end{proof}

\begin{corollary}[The forced problem is a control problem]\label{cor:control}
For fixed $T$, the map \eqref{eq:force} is a bijection between smooth divergence-free paths on
$[0,T)$ with $v(0)=u_0$ and the smooth forces for which \eqref{eq:NS} has a smooth solution
with that datum. Consequently:
\begin{enumerate}[label=(\roman*),leftmargin=*]
\item the \emph{existence} question for the forced problem is trivial: every smooth force that
      arises this way has a smooth solution, by construction;
\item the non-trivial questions are quantitative: which forces are compatible with which
      \emph{a priori} bounds, and how large $f$ must be to produce a given growth of $u$;
\item the unforced problem is the single case $f\equiv0$, i.e.\ the single path satisfying the
      nonlinear constraint \eqref{eq:force} with $f=0$.
\end{enumerate}
\end{corollary}

\begin{proof}
Surjectivity onto paths is Theorem~\ref{thm:prescription}. Injectivity: if two smooth
divergence-free paths $v_1,v_2$ with the same initial datum produce the same $f$, then both
are smooth solutions of the same forced problem, and smooth solutions are unique by the
standard energy argument (the difference $w$ satisfies an equation whose nonlinear terms are
estimated exactly as in Proposition~\ref{prop:uniqueness2d}).
\end{proof}

\begin{corollary}[Manufactured blow-up]\label{cor:manufactured}
Let $T>0$, $a>0$, let $U\in C^\infty_c(\R^3;\R^3)$ be divergence-free and not identically
zero, and put
\begin{equation}\label{eq:selfsimilar}
\lambda(t)=(2a(T-t))^{-1/2},\qquad v(x,t)=\lambda(t)\,U\bigl(\lambda(t)x\bigr).
\end{equation}
Then $v\in C^\infty(\R^3\times[0,T))$, $\Div v=0$, and for every $s>1/2$,
\begin{equation}\label{eq:blowuprate}
\norm{v(t)}_{\dot H^s}=\lambda(t)^{s-1/2}\norm{U}_{\dot H^s}\longrightarrow\infty,
\qquad
\norm{v(t)}_{H^s}\asymp\lambda(t)^{s-1/2}\ (\lambda\ge1),
\qquad
\norm{v(t)}_{\infty}=\lambda(t)\norm{U}_{\infty}\longrightarrow\infty
\end{equation}
as $t\uparrow T$. The force \eqref{eq:force} associated with $v$ has the form
\begin{equation}\label{eq:selfsimilarforce}
f(x,t)=\lambda(t)^3F\bigl(\lambda(t)x\bigr),\qquad
F=a\bigl(U+y\cdot\nabla U\bigr)+(U\cdot\nabla)U-\nu\Delta U+\nabla Q,
\end{equation}
where $Q$ solves $-\Delta Q=\Div((U\cdot\nabla)U)$, and
\begin{equation}\label{eq:forcenorm}
\norm{f(t)}_{\infty}=\lambda(t)^3\norm{F}_{\infty}=(2a(T-t))^{-3/2}\norm{F}_{\infty}.
\end{equation}
For the explicit profile $U$ of Appendix~\ref{app:num}, $\norm{F}_{\infty}>0$; hence for that
profile $\norm{f(t)}_{\infty}\to\infty$ as $t\uparrow T$ and
$\int_0^T\norm{f(t)}_{\infty}\dd t=\infty$, while $v$ is a smooth solution of the forced
problem on $[0,T)$.
\end{corollary}

\begin{proof}
Divergence-free and smoothness are immediate from the definition. The norm identities follow
from \eqref{eq:scaling}: for $v=\lambda U(\lambda\cdot)$ one has
$\norm{v}_{\dot H^s}=\lambda^{s-1/2}\norm{U}_{\dot H^s}$ and
$\norm{v}_\infty=\lambda\norm{U}_\infty$, and $\lambda(t)\to\infty$.

For \eqref{eq:selfsimilarforce}, compute with $y=\lambda x$ and $\lambda=\lambda(t)$, using
$\dot\lambda=a\lambda^3$ (which follows from $\lambda^2=1/(2a(T-t))$):
\[
\partial_tv=\dot\lambda\bigl(U+y\cdot\nabla U\bigr)=a\lambda^3\bigl(U+y\cdot\nabla U\bigr),
\quad
(v\cdot\nabla)v=\lambda^3(U\cdot\nabla)U,\quad
\Delta v=\lambda^3\Delta U,\quad
\nabla q=\lambda^3\nabla Q,
\]
where $q=\lambda^2Q$. Substituting into \eqref{eq:force} gives \eqref{eq:selfsimilarforce}.
Taking $L^\infty$ norms gives \eqref{eq:forcenorm}. The positivity of $\norm{F}_\infty$ for
the profile of Appendix~\ref{app:num} is checked there numerically, with the computation
reproduced in the repository.
\end{proof}

\begin{remark}[The point of the corollary]\label{rem:manufactured}
Corollary~\ref{cor:manufactured} produces a finite-time blow-up of a smooth solution at any
prescribed time $T$, in every $H^s$ with $s>1/2$, for a smooth force. It is therefore a
counterexample to the statement ``smooth force implies global smooth solution'' in the
strongest possible form. What it is \emph{not} is a counterexample to the unforced problem:
the force \eqref{eq:selfsimilarforce} is not merely non-zero, it degenerates at the rate
$(T-t)^{-3/2}$, and Theorem~\ref{thm:price} shows that this rate is forced by the scaling of
the equations.
\end{remark}

\section{The price of a self-similar blow-up}\label{sec:price}

Corollary~\ref{cor:manufactured} shows that blow-up can be bought with a force that
degenerates like $(T-t)^{-3/2}$. The next theorem shows that this is not an artefact of the
construction: for an \emph{exactly} self-similar blow-up the rate is forced, with equality.

\begin{theorem}[Scaling law for the force]\label{thm:price}
Let $\nu>0$, $T\in(0,\infty)$, $a>0$, and let
$\lambda(t)=(2a(T-t))^{-1/2}$. Suppose that $(u,p,f)$ is a classical solution of \eqref{eq:NS}
on $\R^3\times(0,T)$ which is exactly self-similar, i.e.\ there are
$U,P\in C^\infty(\R^3)$ with
\begin{equation}\label{eq:selfsimsol}
u(x,t)=\lambda(t)U\bigl(\lambda(t)x\bigr),
\qquad
p(x,t)=\lambda(t)^2P\bigl(\lambda(t)x\bigr).
\end{equation}
Then there is $F\in C^\infty(\R^3;\R^3)$ such that
\begin{equation}\label{eq:selfsimforce2}
f(x,t)=\lambda(t)^3F\bigl(\lambda(t)x\bigr)
\end{equation}
for all $(x,t)\in\R^3\times(0,T)$, namely
\begin{equation}\label{eq:Fformula}
F(y)=a\bigl(U+y\cdot\nabla U\bigr)+(U\cdot\nabla)U-\nu\Delta U+\nabla P .
\end{equation}
Consequently, for every $t\in(0,T)$,
\begin{equation}\label{eq:exactrate}
\norm{f(t)}_{\infty}=(2a(T-t))^{-3/2}\norm{F}_{\infty},
\qquad
\norm{f(t)}_{q}=(2a(T-t))^{-(3-3/q)}\norm{F}_{q},
\end{equation}
and $\int_0^T\norm{f(t)}_{\infty}\dd t=\infty$ whenever $F\not\equiv0$.
\end{theorem}

\begin{proof}
By Theorem~\ref{thm:prescription}, $f=\partial_tu+(u\cdot\nabla)u-\nu\Delta u+\nabla p$. Put
$y=\lambda(t)x$ and note that $\dot\lambda=a\lambda^3$, which follows from
$\lambda^{-2}=2a(T-t)$ by differentiation: $-2\lambda^{-3}\dot\lambda=-2a$, i.e.\
$\dot\lambda=a\lambda^3$. The chain rule gives, exactly as in the proof of
Corollary~\ref{cor:manufactured},
\[
\partial_tu=a\lambda^3\bigl(U+y\cdot\nabla U\bigr),\qquad
(u\cdot\nabla)u=\lambda^3(U\cdot\nabla)U,\qquad
\Delta u=\lambda^3\Delta U,\qquad
\nabla p=\lambda^3\nabla P .
\]
Adding, $f=\lambda^3F(y)$ with $F$ given by \eqref{eq:Fformula}. Taking $L^q$ norms and using
$\norm{f(t)}_{q}=\lambda^{3-3/q}\norm{F}_{q}$ (which is \eqref{eq:forcescaling} with $F$ in
place of $f$) gives \eqref{eq:exactrate}. Finally, if $F\not\equiv0$, then
$\norm{F}_\infty>0$ and $\int_0^T\lambda^3\dd t=\infty$ because $\lambda^3=(2a)^{-3/2}(T-t)^{-3/2}$
is not integrable at $T$.
\end{proof}

\begin{corollary}[A bounded force cannot support exact self-similarity]\label{cor:price}
In the setting of Theorem~\ref{thm:price}, if $\sup_{0<t<T}\norm{f(t)}_{\infty}<\infty$, then
$F\equiv0$ and $f\equiv0$ on $\R^3\times(0,T)$. Equivalently: a force which is bounded in
$L^\infty$ and not identically zero cannot accompany an exactly self-similar blow-up with
self-similar pressure.
\end{corollary}

\begin{proof}
By \eqref{eq:exactrate}, $\lambda(t)^3\norm{F}_\infty\le M$ for all $t<T$. Since
$\lambda(t)\to\infty$, this forces $\norm{F}_\infty=0$, so $F\equiv0$ and hence $f\equiv0$.
\end{proof}

\begin{corollary}[Discrete self-similarity]\label{cor:discrete}
Suppose that along a sequence $t_k\uparrow T$ the solution is self-similar with a fixed profile
in the sense that $u(x,t_k)=\lambda_kU(\lambda_kx)$ with $\lambda_k\to\infty$. Then
\[
\limsup_{k\to\infty}(T-t_k)^{3/2}\norm{f(t_k)}_{\infty}
=(2a)^{-3/2}\norm{F}_{\infty}.
\]
In particular, a force bounded in $L^\infty$ forces $F\equiv0$ along the sequence.
\end{corollary}

\begin{proof}
Apply Theorem~\ref{thm:price} at each $t_k$ and take the limit.
\end{proof}

\begin{remark}[Reading of the exponent]\label{rem:exponent}
The exponent $3/2$ is nothing but the scaling dimension of the force: $f_\lambda=\lambda^3f$
by Proposition~\ref{prop:scaling}, and a self-similar blow-up with time scale
$\tau=T-t$ has spatial scale $\tau^{1/2}$, so the force at scale $\tau^{1/2}$ is
$\tau^{-3/2}$ times its value at scale one. Equivalently,
\[
\norm{f}_{L^1(0,T;L^\infty)}=\infty
\quad\text{and}\quad
\norm{f}_{L^{2/3}(0,T;L^\infty)}=\infty,
\]
while $\norm{f}_{L^p(0,T;L^\infty)}<\infty$ for every $p<2/3$. Theorem~\ref{thm:price} is a
statement about the \emph{exact} self-similar ansatz; it does not exclude Type~II or
non-self-similar scenarios, which is precisely where the difficulty of the unforced problem
lies.
\end{remark}

\begin{remark}[Relation to \cite{CIV2026}]\label{rem:CIVrelation}
Theorem~\ref{thm:price} is a scaling statement and requires exact self-similarity.
Theorem~\ref{thm:CIV} below, quoted from \cite{CIV2026}, is an analytic statement: within a structural
class, a force that is real-analytic in space (uniformly in time) makes the solution regular,
so the force in the 2026 construction must fail to be analytic near the singular point. The
two statements point at the same phenomenon from different sides: the force must degenerate,
and Theorem~\ref{thm:price} quantifies the rate in the self-similar case.
\end{remark}

\section{The 2026 construction: what is proved and what is blocked}\label{sec:2026}

We now specialise to the announced forced blow-up. Our aim is not to referee the manuscript
--- we have not verified it --- but to state precisely which quoted theorem blocks which
extension, so that the logical situation is unambiguous.

\begin{theorem}[Constantin--Ignatova--Vicol {\rm\cite{CIV2026}}]\label{thm:CIV}
Let $(u,p,f)$ be a suitable weak solution of \eqref{eq:NS} on a parabolic cylinder
$Q\subset\R^3\times(0,T)$ with $f$ real-analytic in the space variables, locally uniformly in
time. Assume
\begin{enumerate}[label=(\roman*),leftmargin=*]
\item the angular mean of $u$ satisfies anisotropic Type~II bounds, and
\item $u$ is exactly axisymmetric in a collapsing core region,
\end{enumerate}
in the precise sense of \cite{CIV2026}. Then $u$ is regular at the putative singular point.
Consequently, in the construction of the 2026 manuscript, and in any construction with
properties (i)--(ii) whose force remains bounded in $C^2$ up to the singular time, the force
can neither vanish identically near the singular point nor be real-analytic in the space
variables, locally uniformly in time.
\end{theorem}

Theorem~\ref{thm:CIV} is quoted, not reproved: the proof rests on unique continuation for the
linearised Stokes system combined with a zooming-in argument at the putative singularity, and
is beyond the scope of this paper. Its companion \cite{IKX2026} replaces spatial analyticity
of $f$ by rotational invariance of $\curl f$.

\begin{corollary}[What the construction must do]\label{cor:CIV}
Assume the 2026 construction satisfies (i)--(ii) of Theorem~\ref{thm:CIV}. Then its force
satisfies
\begin{enumerate}[label=(\roman*),leftmargin=*]
\item $f\not\equiv0$ in every neighbourhood of the singular point (in particular, the argument
      cannot be localised to produce a solution of the unforced problem);
\item $f$ is not real-analytic in $x$ uniformly in $t$ near the singular point, and is not
      bounded in $C^2$ there.
\end{enumerate}
In particular no argument of the form ``the force is supported away from the singularity,
hence the singular behaviour is that of the unforced equation'' can work within this class.
\end{corollary}

\begin{proof}
Immediate from Theorem~\ref{thm:CIV}: $f\equiv0$ is real-analytic and bounded in $C^2$, so
either alternative would make the solution regular.
\end{proof}

\begin{remark}[Two different mechanisms of obstruction]\label{rem:twomech}
Corollary~\ref{cor:CIV} and Theorem~\ref{thm:price} are logically independent.
Theorem~\ref{thm:price} is quantitative and applies to exact self-similarity of any solution,
forced or not, and says the force must blow up like $(T-t)^{-3/2}$ in that scenario.
Corollary~\ref{cor:CIV} is qualitative and applies to a structural class of solutions, and
says the force cannot be analytic or $C^2$-bounded there. Neither implies the other, and
neither settles the unforced problem: the case $f\equiv0$ evades both (it is analytic, so
Theorem~\ref{thm:CIV} would give regularity \emph{if} the structural hypotheses (i)--(ii) were
known to hold for an unforced singular solution; they are not).
\end{remark}

\begin{remark}[On the ``wrong problem'' objection]\label{rem:wrong}
The official statement \cite{Fefferman2000} permits a prescribed smooth force; the discussion
in \cite{SciAm2026} calls this option (C). Theorem~\ref{thm:prescription} shows that option
(C), read literally as ``does a smooth solution exist for smooth data and smooth force'', has
an affirmative answer for a trivial reason. The content of a result for option (C) is
therefore entirely in the \emph{class} of admissible forces: if the class is
$C^\infty(\R^3\times[0,T])$ with all derivatives bounded, a blow-up is a strong statement; if
the class only requires smoothness on $[0,T)$, Corollary~\ref{cor:manufactured} produces
blow-ups with $\norm{f(t)}_\infty\sim(T-t)^{-3/2}$ and Corollary~\ref{cor:CIV} says that this
degeneration is unavoidable in the announced class. This is the precise sense in which the
2026 result answers a different question from the unforced problem; it is not a rhetorical
objection but a statement about quantifiers.
\end{remark}

\section{Machine-checked proofs and the autoformalisation gap}\label{sec:lean}

A separate issue concerns the evidential value of the Lean formalisation that accompanied the
announcement. The following is quoted.

\begin{proposition}[Bastounis--Circelli--Hansen {\rm\cite{BCH2026}}]\label{prop:lean}
Semantically faithful autoformalisation of natural-language mathematics into a formal system
has Solvability Complexity Index $\infty$; informally, it is harder than the halting problem.
Consequently no algorithm can certify that a formal proof is a faithful translation of a given
natural-language proof, and there are explicit examples --- including the Navier--Stokes
manuscript of 2026 --- in which the Lean formalisation does not correspond to the
natural-language argument.
\end{proposition}

The practical consequence for the present paper is a constraint on what may be claimed as
evidence:
\begin{enumerate}[label=(\alph*),leftmargin=*]
\item a verified Lean file certifies that \emph{some} formal statement has a proof; it does not
      certify that the formal statement is the theorem stated in the paper;
\item the gap is not a bug that a better translator would remove: by
      Proposition~\ref{prop:lean} it is a hardness result;
\item therefore the status of the 2026 claim must be assessed on the natural-language
      manuscript, which is what the Clay Institute's process does \cite{ClayNS}.
\end{enumerate}
None of this says that the formalisation is wrong, only that it is not, by itself,
verification of the paper.

\section{A checklist for announced proofs}\label{sec:referee}

The following questions are the ones that decided the status of the 2026 claim, and they are
the questions any future claim must answer.

\begin{enumerate}[label=\textbf{Q\arabic*.},leftmargin=*]
\item \textbf{Which formulation?} Unforced or forced; $\R^3$ or $\T^3$; decay or periodicity?
      The answers are not interchangeable, and Theorem~\ref{thm:prescription} shows that the
      forced version is a different problem.
\item \textbf{Which class of forces?} Smooth on $[0,T)$, or smooth and bounded with all
      derivatives on $[0,T]$? Corollary~\ref{cor:manufactured} shows that the first class
      admits blow-ups of size $(T-t)^{-3/2}$ in the force; for the second class no such
      construction is available, and exact self-similarity is excluded by
      Corollary~\ref{cor:price}.
\item \textbf{Is the force allowed to degenerate at the singular time?} If yes, what is the
      rate, and is it compatible with Theorem~\ref{thm:price}?
\item \textbf{Which structural hypotheses are used?} Axisymmetry, analyticity, Type~I or
      Type~II bounds? Each must be either proved from the data or stated as a hypothesis, and
      Theorem~\ref{thm:CIV} shows that some of them are incompatible with a bounded analytic
      force.
\item \textbf{Are constants uniform up to $T$?} A bound that degrades as $t\uparrow T$ cannot
      be used for continuation; see the proofs of Theorems~\ref{thm:PS} and~\ref{thm:BKM}.
\item \textbf{Is the formalisation faithful?} See \S\ref{sec:lean}; a Lean check of a
      translation is not a check of the paper.
\item \textbf{Is the computation reproducible?} For computer-assisted steps, are the code,
      the inputs and the numerical tolerances public? Appendix~\ref{app:num} is an example of
      what this looks like at small scale.
\end{enumerate}

\section{Conclusion}\label{sec:conclusion}

We have not solved the Navier--Stokes Millennium Problem, and nothing above should be read as
an attempt to. What the paper establishes is a sharp separation of statements:
\begin{itemize}[leftmargin=*]
\item For $f\equiv0$ in three dimensions, the known positive results are conditional
      (Theorems~\ref{thm:PS}, \ref{thm:BKM}) or partial (Theorem~\ref{thm:CKN}); the global
      regularity question is open, and the only complete unconditional theory is the
      two-dimensional one (Theorem~\ref{thm:2d}).
\item For free $f$, the problem is a control problem (Theorem~\ref{thm:prescription}): every
      smooth divergence-free path is a solution for a unique force, and blow-up in every $H^s$,
      $s>1/2$, can be manufactured at any prescribed time
      (Corollary~\ref{cor:manufactured}).
\item The price of a self-similar blow-up is a force of size $(T-t)^{-3/2}$
      (Theorem~\ref{thm:price}), and in the structural class of the 2026 construction the force
      cannot be analytic or $C^2$-bounded at the singular time (Theorem~\ref{thm:CIV},
      Corollary~\ref{cor:CIV}).
\end{itemize}
The third item is the precise sense in which the 2026 announcement does not, and cannot by
localisation, settle the unforced problem. Whether the unforced problem is true or false
remains unknown; what is known is that any proof must produce a mechanism for
$\norm{\nabla u}_\infty\to\infty$ that survives the diffusive smoothing at the critical scale,
and no such mechanism has been found in two centuries of work on the equations.

\appendix

\section{Inequalities used}\label{app:ineq}

For completeness we record the standard inequalities invoked above, in the form used.

\begin{enumerate}[label=\textbf{(A.\arabic*)},leftmargin=*]
\item \emph{Gagliardo--Nirenberg (3D).} For $q>3$ and $r=2q/(q-2)$,
\[
\norm{\nabla u}_{L^r(\R^3)}\le C\norm{\nabla u}_{2}^{1-3/q}\norm{\nabla^2u}_{2}^{3/q},
\]
which follows from $\dot H^1\subset L^6$ and interpolation; the exponents satisfy
$1/r=1/2-\theta/3$ with $\theta=3/q$. Used in Theorem~\ref{thm:PS}.

\item \emph{Sobolev embeddings.} In two dimensions $W^{1,p}\subset L^\infty$ and
$W^{2,p}\subset W^{1,\infty}$ for $p>2$; in three dimensions $W^{1,p}\subset L^\infty$ for
$p>3$. Used in Lemma~\ref{lem:biot} and Theorem~\ref{thm:2d}.

\item \emph{Riesz transforms.} The operators $\partial_i\partial_j(-\Delta)^{-1}$ and
$\partial_i(-\Delta)^{-1/2}$ are bounded on $L^p$ for $1<p<\infty$; this is the content of
\eqref{eq:biot} and of the Biot--Savart representation in two dimensions.

\item \emph{Kato--Ponce commutator estimate.} For $s>0$ and Schwartz functions $f,g$,
\[
\norm{[\Lambda^s,f]g}_{2}\le C\bigl(\norm{\nabla f}_{\infty}\norm{g}_{H^{s-1}}
+\norm{g}_{\infty}\norm{\nabla f}_{H^{s-1}}\bigr).
\]
For divergence-free $u$ this gives
$\abs{\ip{[\Lambda^s,u\cdot\nabla]u}{\Lambda^su}}\le C\norm{\nabla u}_\infty
\norm{u}_{H^s}^2$, used in Theorem~\ref{thm:BKM}.

\item \emph{Kato's inequality.} For $\omega\in C^2$ and $p\ge2$,
$\Delta\abs{\omega}^p\ge p\abs{\omega}^{p-2}\omega\Delta\omega$ in the sense of
distributions. Used in Lemma~\ref{lem:vorticity}.

\item \emph{Ladyzhenskaya inequality (2D).} $\norm{w}_{4}^2\le C\norm{w}\norm{\nabla w}$
for $w\in H^1(\T^2)$ with zero mean. Used in Proposition~\ref{prop:uniqueness2d}.

\item \emph{Aubin--Lions lemma.} A sequence bounded in $L^2(0,T;H^1)$ with time derivatives
bounded in $L^2(0,T;H^{-1})$ has a strongly convergent subsequence in $L^2(0,T;L^2_{\rm loc})$.
Used in Theorem~\ref{thm:lerayhopf}.

\item \emph{Gronwall's lemma.} If $y' \le a y+b$ with $a,b\in L^1(0,T)$, then
$y(t)\le e^{\int_0^ta}(y(0)+\int_0^tb)$. Used throughout.
\end{enumerate}

\section{Numerical verification}\label{app:num}

The identities of \S\S\ref{sec:scaling}--\ref{sec:price} were checked numerically with the
script \texttt{verify.py} in the repository (plain Python with \texttt{numpy}; no other
dependencies). Every field is a trigonometric polynomial or a rapidly decaying Gaussian on a
large periodic box, so the checks are exact up to round-off. The script runs in a few seconds.

\paragraph{T1 (prescription theorem, Theorem~\ref{thm:prescription}).}
A time-dependent divergence-free path $v(t)=a(t)U_1+b(t)U_2$ with $U_j=\curl A_j$ trigonometric
is inserted into \eqref{eq:force}, the pressure being obtained from \eqref{eq:pressure} by an
FFT solve. Results:
\[
\max\abs{\Div v}=2.1\cdot10^{-13},\qquad
\max\abs{\text{residual}}=0,\qquad
\max\abs{\Div f}=1.5\cdot10^{-10},\qquad
\sup\abs{v}=20.0 .
\]
The third number is the substantive one: $\Div f=0$ holds only if the sign in
\eqref{eq:pressure} is correct, so this checks the pressure convention rather than merely
recomputing the residual.

\paragraph{T2 (scaling, Proposition~\ref{prop:scaling}).}
For arbitrary smooth $(u,p,f)$ (not a solution), the residual
$R(u,p,f)=(u\cdot\nabla)u-\nu\Delta u+\nabla p-f$ is computed, then the scaled triple is
computed by exact spectral dilation ($\hat u_\lambda(k)=\lambda\hat u(k/\lambda)$, etc.), and
the two are compared:
\[
\sup\abs{R(u_\lambda,p_\lambda,f_\lambda)-\lambda^3R(u,p,f)(\lambda\cdot)}
=4.5\cdot10^{-12}\ (\lambda=2),\quad
5.7\cdot10^{-12}\ (\lambda=3),
\]
against $\sup\abs{R}=1.03\cdot10^{2}$, i.e.\ a relative error of $4\cdot10^{-14}$. In
particular the exponent $3$ in $f_\lambda=\lambda^3f$ is confirmed.

\paragraph{T3 (price law, Theorem~\ref{thm:price}).}
With the Gaussian profile $\psi=e^{-\abs{y}^2/2}$, $U=\curl(\psi e_z)$, $a=1$, $\nu=1$, a
$96^3$ grid on $[-16,16]^3$:
\[
\sup\abs{\Div U}=8.9\cdot10^{-16},\qquad
\norm{F}_{\infty}=3.262263 .
\]
The chain rule $\partial_tv=\dot\lambda(U+y\cdot\nabla U)$ at $\lambda=1$ (where
$\dot\lambda/\lambda^3=1$) is confirmed to $2.6\cdot10^{-6}$, the size of the $O(h^2)$
finite-difference error with $h=10^{-3}$. Finally the price law
\begin{equation}\label{eq:numerics}
(T-t)^{3/2}\norm{f(t)}_{\infty}=(2a)^{-3/2}\norm{F}_{\infty}=1.153384
\end{equation}
is independent of $t$ to all printed digits for $t=0.25,\,0.5,\,0.875$ (i.e.\
$\lambda=0.8165,\,1,\,2$), which is the numerical form of \eqref{eq:exactrate}.

\paragraph{T4 (vorticity cancellation, Lemma~\ref{lem:vorticity}).}
For a random divergence-free trigonometric $u$ on $\T^2$ with $\omega=\curl u$,
\[
\abs{\int_{\T^2}(u\cdot\nabla)\omega\,\omega\dd x}=6.6\cdot10^{-15},
\]
against a normalisation $\norm{u}_\infty\norm{\nabla\omega}_2\norm{\omega}_2=1.3\cdot10^{4}$,
a relative size of $5\cdot10^{-19}$. This is the identity that kills the transport term in
\eqref{eq:vortbound} and hence is the reason two dimensions are tractable.

\paragraph{What the numbers do and do not show.}
They confirm that the algebra of \S\S\ref{sec:scaling}--\ref{sec:price} is implemented
correctly and that the exponents $3$ and $3/2$ are the right ones. They say nothing about the
existence or non-existence of an unforced blow-up, and they are not a substitute for the
proofs: T1--T4 verify identities, not theorems whose content is an \emph{a priori} bound.

\begin{thebibliography}{99}

\bibitem{Fefferman2000}
C.~L. Fefferman,
\emph{Existence and smoothness of the Navier--Stokes equation},
Clay Mathematics Institute Millennium Prize Problems, 2000.

\bibitem{Leray1934}
J.~Leray,
\emph{Sur le mouvement d'un liquide visqueux emplissant l'espace},
Acta Math. \textbf{63} (1934), 193--248.

\bibitem{Hopf1951}
E.~Hopf,
\emph{\"Uber die Anfangswertaufgabe f\"ur die hydrodynamischen Grundgleichungen},
Math. Nachr. \textbf{4} (1951), 213--231.

\bibitem{Kato1984}
T.~Kato,
\emph{Strong $L^p$-solutions of the Navier--Stokes equation in $\R^m$, with applications to
weak solutions},
Math. Z. \textbf{187} (1984), 471--480.

\bibitem{RRST2001}
J.~C. Robinson, J.~L. Rodrigo, W.~Sadowski,
\emph{The Three-Dimensional Navier--Stokes Equations: Classical Theory},
Cambridge Studies in Advanced Mathematics 157, Cambridge University Press, 2016.

\bibitem{Lemarie2016}
P.~G. Lemari\'e-Rieusset,
\emph{The Navier--Stokes Problem in the 21st Century}, 2nd ed.,
CRC Press, 2016.

\bibitem{Prodi1959}
G.~Prodi,
\emph{Un teorema di unicit\`a per le equazioni di Navier--Stokes},
Ann. Mat. Pura Appl. \textbf{48} (1959), 173--182.

\bibitem{Serrin1962}
J.~Serrin,
\emph{On the interior regularity of weak solutions of the Navier--Stokes equations},
Arch. Rational Mech. Anal. \textbf{9} (1962), 187--195.

\bibitem{BKM1984}
J.~T. Beale, T.~Kato, A.~Majda,
\emph{Remarks on the breakdown of smooth solutions for the $3$-D Euler equations},
Comm. Math. Phys. \textbf{94} (1984), 61--66.

\bibitem{Ladyzhenskaya1969}
O.~A. Ladyzhenskaya,
\emph{The Mathematical Theory of Viscous Incompressible Flow}, 2nd ed.,
Gordon and Breach, 1969.

\bibitem{CKN1982}
L.~Caffarelli, R.~Kohn, L.~Nirenberg,
\emph{Partial regularity of suitable weak solutions of the Navier--Stokes equations},
Comm. Pure Appl. Math. \textbf{35} (1982), 771--831.

\bibitem{GigaSohr1991}
Y.~Giga, H.~Sohr,
\emph{Abstract $L^p$ estimates for the Cauchy problem with applications to the Navier--Stokes
equations in exterior domains},
J. Funct. Anal. \textbf{102} (1991), 72--94.

\bibitem{CIV2026}
P.~Constantin, M.~Ignatova, V.~Vicol,
\emph{Regularity of asymptotically axisymmetric solutions to the 3D Navier--Stokes equations
with analytic forcing}, arXiv:2609.20803, 2026.

\bibitem{IKX2026}
B.~Ingimarson, I.~Kukavica, Q.~Xu,
\emph{Regularity of asymptotically axisymmetric Navier--Stokes flows with rotationally
invariant curl of the force}, arXiv:2610.25531, 2026.

\bibitem{BCH2026}
A.~Bastounis, F.~Circelli, A.~C. Hansen,
\emph{Navier--Stokes lost in translation: why Lean verification of AI autoformalisation does
not guarantee correct natural language proofs}, arXiv:2610.08144, 2026.

\bibitem{OpenAI2026}
OpenAI, \emph{On the Navier--Stokes Millennium Prize problem}, manuscript and Lean
formalisation, September 2026.

\bibitem{OpenAIrepo2026}
OpenAI, \emph{Mathematics manuscript collection} (721 manuscripts, 373 result families),
\url{https://github.com/openai/math}, accessed October 2026.

\bibitem{RW2026}
A.~Gallegos,
\emph{OpenAI withdraws three preprints a day after releasing 722 manuscripts on unsolved math
problems}, Retraction Watch, 8 October 2026.

\bibitem{SciAm2026}
J.~Howlett,
\emph{Did OpenAI solve the wrong Navier--Stokes problem?}, Scientific American, 21 September
2026.

\bibitem{ClayNS}
Clay Mathematics Institute,
\emph{Navier--Stokes equation} (status: active),
\url{https://www.claymath.org/millennium/navier-stokes-equation/}, accessed November 2026.

\bibitem{WikipediaNS}
Wikipedia, \emph{Navier--Stokes existence and smoothness}, section on the 2026 announcement,
accessed November 2026.

\end{thebibliography}

\end{document}















